\subsection{Coverage and the frequency law}
\label{sec:coverage}
Section~\ref{sec:forking-factors} found that low-frequency contexts carry most of the gap. A finite training corpus leaves some plausible continuations of each context unobserved. Let $f(F)$ be the fraction of validation positions whose context occurs $F\ge1$ times in training but whose true continuation never followed it there---\emph{seen context, novel continuation}. Averaged over a frequency bin, the gap splits exactly into a novel and a seen term,
\begin{equation}
\begin{aligned}
\gap_t(F)&=f(F)D_t(F)+\big(1-f(F)\big)S_t(F),\\
D_t(F)&=\ell^{\mathrm{val}}_{\mathrm{novel},t}(F)-\ell^{\mathrm{tr}}_t(F),\qquad
S_t(F)=\ell^{\mathrm{val}}_{\mathrm{seen},t}(F)-\ell^{\mathrm{tr}}_t(F).
\end{aligned}
\label{eq:gap-decomp}
\end{equation}
Here $\ell^{\mathrm{tr}}_t$ is the current-batch training loss and the validation means use the fixed validation pool at the same logged step. These are means within a frequency bin, not losses matched context by context. Coverage sets the weights; $D_t$ and $S_t$ describe how the model responds. The seen term need not be small.

\paragraph{From coverage to the frequency law.}
If $D_t$ and $S_t$ were constant across $F$, the decomposition would be affine in $f$. We test the corresponding two-parameter approximation directly,
\begin{equation}
\gap_t(F)\approx a_t f(F)+b_t.
\label{eq:gap_j}
\end{equation}
For the both-table model, 18 frequency bins spanning $F=1$--$1000$ lie close to a line in the directly measured novel fraction (Figure~\ref{fig:coverage}b,c). For fits to the three-seed mean at steps 670, 1000, 1340 and 2000, relative rms residuals are 3--8\% for bigram and 3--11\% for trigram contexts. The intercept matters: at step 2000 the bigram fit has $b_t=1.89$ nats; forcing the line through zero increases rms error from 0.20 to 0.76 nats. This approximate affine relation summarizes the frequency dependence, but does not establish constant component losses: a fitted $b_t$ is not in general the measured $S_t(F)$ (Appendix~\ref{app:coverage-tests}).

\paragraph{Coverage from corpus counts.}
For a context with independent continuation draws, the expected singleton fraction after $F+1$ observations equals the expected missing mass after $F$. Pooling heterogeneous contexts additionally requires a model for their occurrence counts. Under Poisson counts and matched train/validation distributions, the token-weighted population novel fraction satisfies
\begin{equation}
f_\star(F)=\frac{\mathbb E N_1(F+1)}{(F+1)\mathbb E n(F+1)},\qquad
\widehat f_{\mathrm{GT}}(F)=\frac{N_1(F+1)}{(F+1)n(F+1)},
\label{eq:gt}
\end{equation}
where $n(k)$ counts contexts observed $k$ times and $N_1(k)$ their singleton context--continuation types. The first equality is a Good--Turing expectation identity; the second defines its empirical estimator \citep{good1953population,gale1995good}. The derivation and assumptions are in Appendix~\ref{app:gt-identity}. On our corpus the uncalibrated estimate tracks held-out novelty but underestimates it: its ratio to the direct count is 0.870--0.885 for bigram and 0.857--0.937 for trigram contexts (Figure~\ref{fig:coverage}a). No validation-fitted correction is applied.

\begin{figure}[!t]
\centering
\includegraphics[width=\linewidth]{fig_coverage_linear.pdf}
\caption{Coverage and the gap. (a) Novel fraction predicted from training counts against directly measured held-out novelty; the dashed diagonal is equality. (b,c) Online gap against the directly measured novel fraction, with affine fits at four training steps. Each point is a frequency bin; bars in (b,c) show the standard deviation across three seeds. The no-gram control is shown at step 1000. All panels use linear axes; protocol and residuals are in Appendix~\ref{app:coverage-tests}.}
\label{fig:coverage}
\end{figure}

Coverage also varies among contexts with the same count. At exact $F=4$, splitting bigram contexts by their training singleton fraction gives novel fractions 0.25 versus 0.80 and diagnostic gaps 3.23 versus 6.01 nats in the extreme strata (seed 42, step 1000; fixed-probe diagnostic in Appendix~\ref{app:coverage-tests}). This supports a role for continuation statistics beyond count alone, without isolating coverage from conditional entropy or optimization. Conditional heavy tails explain why missing mass can decay slowly, but do not fix a cross-context gap exponent. Coverage describes which continuations lack evidence; the replay dynamics below determine the response to that missing evidence.
